% index_prompt_generator.tex
\documentclass[12pt]{article}
\usepackage[utf8]{inputenc}
\usepackage[paper=a4paper,margin=2.5cm]{geometry}
\begin{document}

\section*{אינדקס prompt\_generator.tex}
\textbf{מטרה}:
להגדיר את מבנה הפרומט המחקרי שנוצר אוטומטית.

\subsection*{סעיפי הקובץ}
\begin{itemize}
  \item Title שם המחקר.
  \item Objective מטרת המחקר בשורה אחת.
  \item Scope גבולות המחקר.
  \item Research Questions רשימת שאלות מחקר.
  \item Methodology פירוט שיטות.
  \item Data Sources רשימת מקורות נתונים.
  \item Inclusion Criteria קריטריוני הכללה.
  \item Exclusion Criteria קריטריוני שלילה.
  \item Extraction Fields שדות לחילוץ.
  \item Deliverables תוצרים צפויים.
  \item Constraints מגבלות אתיות וטכניות.
  \item Evaluation Metrics מדדי הערכה.
  \item Tone and Style סגנון הכתיבה.
  \item Length and Format פורמט ואורך הפלט.
\end{itemize}

\subsection*{הפניות}
הסתכלו ב-\textbf{index\_parameters.tex} להגדרות מפורטות של כל פרמטר.
\end{document}
